\documentclass[a4paper,11pt]{article}


% Please do not change the packages and commands listed here. You can of course add your favourite packages and commands.

%%% Pakete %%%

\usepackage{graphicx}
\usepackage{amsmath}
\usepackage[utf8]{inputenc}
\usepackage{enumerate}
\usepackage{amsthm}
\usepackage{amssymb}
\usepackage{paralist}
\usepackage{dsfont}
\usepackage{booktabs}
\usepackage{tabularx}
\usepackage{hyperref}
\usepackage{mathtools}
\usepackage{verbatim}
\usepackage{ dsfont }
\usepackage{ xcolor }
\usepackage{environ}
\usepackage{lipsum}


%%% Additional Commands %%%

%% These commands might mak it easier for you to write formal statements.

\newcommand{\N}{\ensuremath{\mathds{N}}} % natural numbers
\newcommand{\Z}{\ensuremath{\mathds{Z}}} % integers
\newcommand{\Q}{\ensuremath{\mathds{Q}}} % rational numbers
\newcommand{\R}{\ensuremath{\mathds{R}}} % real numbers

\newcommand{\states}{\ensuremath{\mathcal{S}}} % state space
\newcommand{\actions}{\ensuremath{\mathcal{A}}} % action space
\newcommand{\policy}{\ensuremath{\pi}} % a policy
\newcommand{\stateval}{\ensuremath{v}} % state value function
\newcommand{\actionval}{\ensuremath{q}} % action value function
\newcommand{\optsval}{\ensuremath{v_*}} % state value function
\newcommand{\optaval}{\ensuremath{q_*}} % action value function
\newcommand{\reward}{\ensuremath{\mathcal{R}}} % reward function
\newcommand{\transition}{\ensuremath{\mathcal{P}}} % transition function
\newcommand{\discount}{\ensuremath{\gamma}} % discount factor

% Wrap AI-assisted text in this to highlight it blue (required by course policy)
\newcommand{\aitxt}[1]{\textcolor{blue}{#1}}

\newtheorem{lemma}{Lemma}
\newtheorem{example}{Example}
\theoremstyle{definition}
\newtheorem{prob}{Exercise}
\newtheorem{pcode}{Pseudocode}




% -----------------------------
\begin{document}

%%% TITLE PAGE %%%
%   - Please only modify your name(s) and do not change anything else! %

\begin{titlepage}
\pagestyle{empty}
  \begin{minipage}[t]{0.6\textwidth}
    \begin{flushleft}
      \bf
      Reinforcement Learning and Decision Making Under Uncertainty\\
    \end{flushleft}
  \end{minipage}\hfill
  \begin{minipage}[t]{0.3\textwidth}
    \begin{flushright}
      \bf Spring 2022\\
      Viola, Meier
      Romano Brentani
      Isabelle Bustamante
    \end{flushright}
  \end{minipage}

	\medskip

  \begin{center}
    {\Large\bf Project Report}

  {\bf Submission Deadline: 03 June 2022, 14:00}
    \Large
    \vspace{2cm}

    %%% ENTER THE NAMES OF GROUP MEMBERS HERE %%%
    %
    {Group:}\\
    {NAME 1}\\
    {NAME 2}

    \vspace{2cm}
    \begin{tabular}[b]{|l|r|}
    \hline
    Introduction, Related Work, Preliminaries &\qquad\qquad/\;\;5\\\hline
    Structure, Notation, Formal correctness &\qquad\qquad/\;\;5\\\hline
    Main Content and Explanations &\qquad\qquad/30\\\hline
    Discussion \& Conclusion&\qquad\qquad/10\\\hline
    $\Sigma$&\qquad\qquad/50\\\hline
    \end{tabular}
  \end{center}
\end{titlepage}
%%% END TITLE PAGE %%%




\newpage
\pagestyle{plain}
\pagenumbering{arabic}


\section{Introduction}
Warehouse automation is an increasingly active area of research and industrial development, with a growing demand for fast, accurate and efficient order fufilment. Inefficient routing increases operational costs and reduces throughput. As warehouses grow in size, finding efficient strategies manually becomes impractical.

We consider a warehouse robot that must pick up packages and deliver them to storage locations in the right order, while avoiding obstacles. This problem is modeled as a Markov Decision Process (MDP) \cite{sutton2018}. The state variable captures the robot's position, the remaining packages, and the locations of obstacles. Possible actions include movement steps and package pickup or delivering packages. The reward system incentivises fast and successful deliveries, while penalising collisions with walls. We implement and evaluate five approaches on a discrete grid warehouse: Q-Learning, SARSA, Value Iteration, Policy Iteration, and Deep Q-Network (DQN) \cite{mnih2015}.
\subsection*{Related Work}
RL has been applied to warehouse robotics across a range of problem settings and scales \cite{punukollu2019integrating, peyas2021warehouse, liu2025drmarl}. A common observation is that warehouse environments are well-suited for RL because they are dynamic and unpredictable, which renders hand-crafted rules insufficient \cite{peyas2021warehouse}.

At the algorithm level, prior work has explored DQN, PPO, and Advantage Actor-Critic for tasks such as route planning, object retrieval, and resource allocation \cite{punukollu2019integrating}. Peyas et al.\ \cite{peyas2021warehouse} use both Deep RL for navigation and Q-learning for product placement in a 2D grid, which is the same action space and environment structure we adopt. At industrial scale, Liu et al.\ \cite{liu2025drmarl} apply multi-agent RL to package sorting in Amazon warehouses and achieve an 80\% reduction in package recirculation under uncertain conditions.

Two challenges appear across this work. Agents require many training episodes before convergence \cite{punukollu2019integrating}. Second, state spaces grow quickly as the problem scales: Peyas et al.\ \cite{peyas2021warehouse} acknowledge that extension to higher-dimensional, real-world environments remains unsolved. In our setting, each additional package doubles the number of states to experiment with higher-dimensional environments. We also test robustness by the addition of dynamic obstacles during execution.

\section{Preliminaries}

\subsection{MDP Formulation}

The warehouse problem is formalised as a tuple $(\states, \actions, \transition, \reward, \discount)$.

\paragraph{State space $\states$.}
Each state is a tuple $s = (\text{pos},\, c,\, \mathbf{d})$, where:
\begin{itemize}
    \item $\text{pos} = (r, c) \in \mathcal{G}$ is the robot's position, with $\mathcal{G}$ the set of passable cells in the grid,
    \item $c \in \{\emptyset, 0, 1, \ldots, N{-}1\}$ is the package currently carried ($\emptyset$ if none),
    \item $\mathbf{d} \in \{0,1\}^N$ is a binary vector indicating which of the $N$ packages have been delivered.
\end{itemize}
States where a delivered package is still being carried are excluded. The total number of states grows as $|\mathcal{G}| \times (N+1) \times 2^N$.

\paragraph{Action space $\actions$.}
There are six actions: $\actions = \{\texttt{up},\, \texttt{down},\, \texttt{left},\, \texttt{right},\, \texttt{pickup},\, \texttt{deliver}\}$. Movement actions are always available; \texttt{pickup} is valid only when the robot stands on an undelivered package and carries nothing; \texttt{deliver} is valid only when the robot carries a package and stands on its designated storage cell.

\paragraph{Transition function $\transition$.}
The environment is deterministic: $\transition(s' \mid s, a) \in \{0, 1\}$ for all $s, a, s'$. Movement to a wall leaves the robot in place. Pickup and deliver actions update $c$ and $\mathbf{d}$ accordingly.

\paragraph{Reward function $\reward$.}
The reward $\reward(s, a)$ is defined as:
\begin{equation}
    \reward(s, a) =
    \begin{cases}
        -1   & \text{valid movement step} \\
        -5   & \text{movement into a wall} \\
        -10  & \text{invalid pickup or deliver} \\
        +9   & \text{successful pickup } (+10 - 1) \\
        +99  & \text{successful delivery } (+100 - 1)
    \end{cases}
\end{equation}

\paragraph{Discount factor $\discount$.}
We use $\discount = 0.95$, which encourages the agent to complete deliveries quickly.

\paragraph{Bellman optimality equation.}
The optimal state-value function $\optsval$ satisfies:
\begin{equation}
    \optsval(s) = \max_{a \in \actions} \left[ \reward(s,a) + \discount \sum_{s' \in \states} \transition(s' \mid s, a)\, \optsval(s') \right]
\end{equation}

\paragraph{Optimal policy.}
The optimal policy selects the action that maximises the expected return in every state:
\begin{equation}
    \policy^*(s) = \arg\max_{a \in \actions} \left[ \reward(s,a) + \discount \sum_{s' \in \states} \transition(s' \mid s, a)\, \optsval(s') \right]
\end{equation}


\section{Environment}

\subsection{Environment Description}
% Grid size, obstacle layout, start state, goal state(s). Include a figure.

\subsection{Implementation Details}
% State encoding, episode termination condition, stochasticity. Link to code repository.


\section{Algorithms}

\subsection{Q-Learning}
% Write the Q-learning update equation and describe the exploration strategy.

\subsection{SARSA}
%

\subsection{Deep Q-Network (DQN)}
% Write the DQN loss function. Mention replay buffer and target network.
\paragraph{State Encoding}

Initial state representation:
\[
s = (\textit{robot\_pos}, \textit{carrying}, \textit{delivered})
\]
where
\[
\textit{robot\_pos} = (row, col),
\]
\[
\textit{carrying} = i \Longleftrightarrow \textit{robot carries package } i,
\]
and
\[
\textit{delivered} = (0,\dots,1,\dots,0)
\]
with a value of 1 at position \(i\) if package \(i\) was successfully delivered.

The carrying state was one-hot encoded and the complete state was passed to the network as a flattened vector.

Results:
\begin{itemize}
    \item rewards remained strictly negative
\end{itemize}

Changed encoding to a grid-based representation.

\paragraph{Grid-based State Encoding}

The MDP was encoded as a multi-channel grid:

\begin{itemize}
    \item channel 0: obstacles (1 if wall at position)
    \item channels 1..N: package start positions
    \item channels N+1..2N: storage locations (0.5 if empty, 1 if delivered)
    \item channel 2N+1: robot position
    \item channels 2N+2..3N+1: carrying state (one channel per package)
\end{itemize}

Notes:
\begin{itemize}
    \item packages and storages may share the same cell
    \item agent may occupy the same position as a package or storage
    \item all channels are concatenated and passed to the network
\end{itemize}

Results remained poor.

Added action masking:
\begin{itemize}
    \item invalid pickup and delivery actions masked
\end{itemize}

Worked temporarily:

\begin{figure}[h!]
  \vskip 0.2in
  \begin{center}
    \centerline{\includegraphics[width=\columnwidth]{report/figures/warehouse_rewards_working.png}}
    \caption{
      Historical locations and number of accepted papers for International
      Machine Learning Conferences (ICML 1993 -- ICML 2008) and International
      Workshops on Machine Learning (ML 1988 -- ML 1992). At the time this
      figure was produced, the number of accepted papers for ICML 2008 was
      unknown and instead estimated.
    }
    \label{fig:deep-q}
  \end{center}
\end{figure}

Training became unstable again using the same setup:

\begin{figure}[h!]
  \vskip 0.2in
  \begin{center}
    \centerline{\includegraphics[width=\columnwidth]{report/figures/warehouse_rewards.png}}
    \caption{
      Historical locations and number of accepted papers for International
      Machine Learning Conferences (ICML 1993 -- ICML 2008) and International
      Workshops on Machine Learning (ML 1988 -- ML 1992). At the time this
      figure was produced, the number of accepted papers for ICML 2008 was
      unknown and instead estimated.
    }
    \label{fig:deep-q}
  \end{center}
\end{figure}

Observations:
\begin{itemize}
    \item successful delivery never occurred
    \item rewards dropped again during training
\end{itemize}

Changed hyperparameters:

\begin{table}[]
\begin{tabular}{|l|l|l|}
\hline
                  & First Training                          & Second Training                         \\ \hline
GAMMA             & 0.99                                    & 0.99                                    \\ \hline
Learning Rate     & 0.001, exp annealed to 0.0001 & 0.001, exp annealed to 0.0001 \\ \hline
EPS\_START        & 0.9                                     & 0.9                                     \\ \hline
EPS\_END          & 0.001                                   & 0.001                                   \\ \hline
EPS\_DECAY        & 1000                                    & 10000                                   \\ \hline
TAU               & 0.005                                   & 0.005                                   \\ \hline
SOFT\_UPDATE      & 0.01                                    & 0.001                                   \\ \hline
NUM\_TRAJECTORIES & 10000                                   & 10000                                   \\ \hline
MAX\_EPS\_LENGTH  & 500                                     & 100                                     \\ \hline
WARMUP            & 1500                                    & 3000                                    \\ \hline
BATCH\_SIZE       & 128                                     & 128                                     \\ \hline
\end{tabular}
\end{table}

Results remained unstable.

\subsection{Hyperparameters}
% Table: learning rate, discount factor, exploration schedule, episodes, and DQN-specific params.


\section{Experiments}

\subsection{Hypotheses}
% List the specific questions the experiments aim to answer.

\subsection{Setup}
% Environment variant(s), number of runs, random seeds, evaluation metric, baselines.

\subsection{Results}
% Learning curves and/or tables. Figure: reward vs.\ episode.

\subsection{Discussion}
% Interpret the results. What explains the observed behaviour?


\section{Robustness}

\subsection{Robustness Analysis}
% Dimensions tested: stochasticity, unseen layouts, reward perturbations, hyperparameter sensitivity.

\subsection{Results}
% Table: performance under each condition. Discuss what the policy tolerates and where it breaks.


\section{Conclusion}

\subsection{Limitations}
% What does the approach not handle well?

\subsection{Future Work}
% Possible extensions to the project.


\section*{Use of External Aids}
% List every book, paper, URL, library, and AI tool used. Highlight AI-assisted text with .


\bibliographystyle{plain}
\bibliography{references}




\end{document}